\section{Introducción al curso y el estado del desarrollo de software en 2025}

\subsection{Por qué se necesita este curso}

En 2025, el desarrollo de software atraviesa un cambio de paradigma profundo. Herramientas de programación con IA como Windsurf, Cursor y Claude Code están cambiando la forma de trabajar de los desarrolladores a una velocidad sin precedentes. Mihail Eric señaló dos señales clave al abrir el curso:

\begin{importantbox}{Idea central: la IA no sustituirá a los desarrolladores}
``You won't be replaced by AI. You'll be replaced by a competent engineer who knows how to use AI.''

La IA en sí misma no sustituirá a los ingenieros de software, pero los ingenieros que sepan usar bien la IA sí sustituirán a los que no sepan usarla. El objetivo de este curso es que los estudiantes sean de los primeros.
\end{importantbox}

La buena noticia es que los desarrolladores de software tienen hoy el mayor potencial de productividad de la historia. Con la ayuda de las herramientas de programación con IA, los ingenieros pueden aprender nuevas pilas tecnológicas y herramientas a una velocidad sin precedentes.

\begin{warningbox}{Esto no es un curso de ``Vibe Coding''}
El curso deja claro que no es un curso que enseñe a generar código con IA sin criterio. El énfasis está en entender los principios de las herramientas de programación con IA y dominar patrones de colaboración human-agent de alta calidad, no en depender de la IA a ciegas.
\end{warningbox}

\subsection{Idea central del curso}

El curso se organiza en torno a las siguientes ideas clave:

\begin{enumerate}
    \item \textbf{Human-Agent Engineering}: se centra en las habilidades que la IA todavía no sustituye: comprensión del negocio, decisiones de arquitectura técnica, convertirse en Tech Lead
    \item \textbf{La capacidad del LLM depende de ti}: un buen context produce buen código; si tú mismo no puedes entender tu base de código, el LLM tampoco puede
    \item \textbf{Leer y revisar mucho código}: aprender a distinguir el código bueno del malo y desarrollar ``criterio'' (taste)
    \item \textbf{Experimentar de forma radical}: por ahora no existe un paradigma maduro de desarrollo de software; todos están explorando
\end{enumerate}

\subsection{Panorama general del curso}

El curso, de 10 semanas, cubre todas las capas de la cadena de herramientas de programación con IA: desde los fundamentos de los LLM, la construcción de Agents de programación, los IDE con IA, las herramientas de terminal, las pruebas y la seguridad, la revisión de código, la construcción automatizada de interfaces de usuario y DevOps, hasta las perspectivas futuras de la ingeniería de software con IA. Cada semana se invita a un ponente destacado del sector.

\subsection{Resumen de la sección}

Esta sección fija el tono del curso: la IA está transformando el desarrollo de software, pero la competitividad central sigue radicando en la comprensión, el juicio y el dominio de las herramientas del propio ingeniero. Lo importante es aprender a colaborar de forma eficiente con los AI Agent, no ser sustituido por la IA.
